\chapter{Attention maps: four ways to fool yourself}
\label{ch:attmaps}

\begin{objectives}
\begin{itemize}
\item Compute what averaging many attention heads does to a heatmap, and why the answer is
  "destroys it".
\item Explain why a high attention weight is not the same as importance.
\item Recognise the artefacts forced by softmax --- attention has to go somewhere.
\item State the orientation question that every board-overlay must answer, and why getting it
  wrong produces a plausible-looking picture.
\item Propose displays that survive all four objections.
\end{itemize}
\end{objectives}

\section{What you are looking at}

Chapter~\ref{ch:attention} produced, for each head $h$ of each layer $l$, a $64\times64$
matrix $A^{(l,h)}$ whose row $i$ is a probability distribution: \emph{how square $i$ divided
its attention over the board}. A heatmap on a chessboard is some reduction of that object down
to 64 numbers.

The reduction is where all the danger is. Recall the compression:
\[
  15 \text{ layers} \times 24 \text{ heads} \times 64 \times 64
  \;=\; 1{,}474{,}560 \text{ numbers} \;\longrightarrow\; 64 .
\]
Any such reduction is a strong claim about what matters. This chapter is about four ways the
claim goes wrong --- and the first one is almost certainly live in your project right now.

\section{Trap 1: averaging heads destroys the signal}

The natural reduction is to average: over heads, over layers, sometimes over both. It is also
close to the worst available choice, because \textbf{heads specialise, and specialists
disagree}.

\begin{worked}{what averaging does to disagreeing heads}
Suppose, over a single square's row, four heads attend as follows across four candidate
squares \emph{(illustrative)}:
\begin{center}
\begin{tabular}{@{}lrrrr@{}}
\toprule
 & sq A & sq B & sq C & sq D \\
\midrule
head 1 (say, attackers) & $0.85$ & $0.05$ & $0.05$ & $0.05$ \\
head 2 (say, defenders) & $0.05$ & $0.85$ & $0.05$ & $0.05$ \\
head 3 (say, pawn structure) & $0.05$ & $0.05$ & $0.85$ & $0.05$ \\
head 4 (say, king safety) & $0.05$ & $0.05$ & $0.05$ & $0.85$ \\
\midrule
\textbf{mean} & $\mathbf{0.25}$ & $\mathbf{0.25}$ & $\mathbf{0.25}$ & $\mathbf{0.25}$ \\
\bottomrule
\end{tabular}
\end{center}
Every head is sharply informative; the average is exactly uniform. Averaging turned four
strong, specific statements into "the network looks at everything equally".
\end{worked}

The general principle: averaging probability distributions moves the result towards the
\emph{most diffuse} of them, and a mean over $H$ diverse distributions has entropy at least as
large as the mean of their entropies. A diffuse heatmap is therefore the \emph{expected}
outcome of averaging specialists, whether or not the network has anything specific to say.

\begin{keyidea}
If your heatmap always looks like a soft cloud around the pieces, that is very likely an
artefact of the averaging, not a fact about the network. \textbf{The test is cheap:} render
two or three individual heads at one layer and see whether they are sharp. If they are, the
averaging is the problem, not the network.
\end{keyidea}

\begin{repobox}[title={in this repository}]
The project notes describe the saliency map as "currently averaged over all layers/heads, so
it is diffuse". This chapter says the diffuseness is predicted by the averaging alone. Before
any deeper work, render single heads --- it costs nothing (the matrices are already extracted)
and it decides whether there is a signal to display at all.
\end{repobox}

\section{Trap 2: attention weight is not importance}

The second failure is subtler, and it is the subject of a genuine research literature
(the "attention is not explanation" debate, from around 2019). Three distinct reasons a
heavily-attended square may not matter:

\begin{enumerate}
\item \textbf{The value vector may be small.} The output is $\sum_j \alpha_{ij}\vv{v}_j$. A
  large $\alpha_{ij}$ multiplied into a near-zero $\vv{v}_j$ contributes nothing. Attention
  says where you looked; the value says whether there was anything to bring back.
\item \textbf{The residual stream may dominate.} Each layer \emph{adds} to what was already
  there (Chapter~\ref{ch:attention}, residual connections). If the attention block's
  contribution is small relative to the existing vector, the whole attention pattern is a
  minor correction regardless of its shape.
\item \textbf{Later layers may discard it.} Information moved at layer 4 can be ignored by
  layers 5--15. This is Chapter~\ref{ch:probes}'s "present but unused", in a different guise.
\end{enumerate}

\begin{keyidea}
\textbf{The honest caption for an attention heatmap is "where this square read from at this
layer", not "what the engine thinks is important".} The second is an importance claim, and
importance claims require the interventions of Chapter~\ref{ch:probes} §13.3.3 --- ablate the
attention edge and show the output changes.
\end{keyidea}

\section{Trap 3: attention has to go somewhere}

Each row is a softmax, so it sums to 1 whether or not the square has anything to attend to. A
square with nothing relevant nearby must still distribute its full unit of attention. This
produces two artefacts you will see and misread:

\begin{itemize}
\item \textbf{Sinks.} Networks commonly learn to dump surplus attention on one or two
  "parking" positions --- in language models this is famously the first token. If one square
  attracts heavy attention across many heads and positions, suspect a sink before concluding
  it is a strategically vital square. \emph{Test: does it attract the same attention in
  unrelated positions?}
\item \textbf{Self-attention.} The diagonal $\alpha_{ii}$ is often the largest entry, meaning
  "keep what I have". It is usually uninformative and, if included in a heatmap's colour
  scale, it flattens everything else. Consider excluding or reporting it separately.
\end{itemize}

\begin{worked}{the same picture from two causes}
Consider a row with a big spike on square X. Two possible explanations, indistinguishable
from the heatmap alone:
\begin{itemize}
\item the network is tracking a real relation to X;
\item nothing here is relevant, and X is where this head parks its surplus.
\end{itemize}
The distinguishing experiment is a \textbf{contrast}: evaluate a near-identical position with
the relation removed (move the piece, block the diagonal) and re-render. A real relation moves
with its cause; a sink does not. Contrast maps are strictly more informative than absolute
maps and cost one extra forward pass.
\end{worked}

\section{Trap 4: orientation, and why it fails silently}

The network's 64 tokens are indexed in some order. Leela's input planes are conventionally
oriented \emph{from the side-to-move's perspective}: when Black is to move, the board is
flipped so that the mover's pieces are always at the bottom. If your display maps token index
$i$ to a fixed board square without accounting for that, then \textbf{every heatmap for a
black-to-move position is reflected}.

\begin{pitfall}
This is the most dangerous bug class in the whole chapter, because the output is not obviously
wrong. A reflected heatmap still looks like a heatmap; it still glows near pieces; it still
tells a story --- just the wrong one, and only in half of all positions. It will pass every
test that does not specifically check orientation, including "looks plausible to an expert".

Your project has already hit this once: the memory note recording that a black-to-move
orientation fix landed as an explicit \texttt{saliency\_absolute(fen)} function, with the older
frame-relative version left in place, is exactly this bug. Note the residual hazard: two
functions, one correct, one not, and nothing in the type system to stop the wrong one being
called.
\end{pitfall}

The general lesson is worth stating beyond this one bug: \textbf{every quantity extracted from
the network carries a frame}, and frames are invisible in the numbers. Values are
side-to-move-relative (Chapter~\ref{ch:uct}). Policy indices are move-encoding-relative. Board
tokens are orientation-relative. Each is a silent convention, and each has produced a
plausible-looking wrong picture in some chess tool. A useful discipline is to make the frame
part of the name of every function that returns one.

\section{Displays that survive}

Given all four traps, here is what can honestly be shown, roughly in order of increasing
strength:

\begin{center}
\begin{tabular}{@{}p{5.1cm}p{8.2cm}@{}}
\toprule
Display & What it honestly claims \\
\midrule
Single head, single layer, one row & "At layer $l$, head $h$, square $e4$ read mainly from
these squares." Weak but true. \\
\addlinespace
Contrast map (this position minus a modified one) & "This part of the pattern depends on the
relation you removed." Much stronger. \\
\addlinespace
Ablation map (zero the edge, measure policy change) & "Removing this changed the engine's
move by this much." An importance claim, earned. \\
\addlinespace
Attention averaged over all heads and layers & Very little. Prefer not to. \\
\bottomrule
\end{tabular}
\end{center}

\begin{keyidea}
The strongest visual claim you can make about a neural network is a \textbf{difference}, not a
level. "This changed when I changed that" is checkable and causal; "this square is bright" is
neither. If the project's visual channel is to carry real weight --- and the design notes argue
it should --- contrast and ablation maps are where its credibility will come from.
\end{keyidea}

\begin{exercises}

\exhead[averaging]{ch14:averaging} Three heads attend to squares $(A,B,C)$ with distributions
$(0.9,0.05,0.05)$, $(0.05,0.9,0.05)$, $(0.45,0.45,0.1)$. Compute the mean. Compute the
entropy of each head and of the mean. Which is more diffuse than \emph{any} individual head,
and does that surprise you?

\exhead[the cheap test]{ch14:cheaptest} Write out, in five steps, the experiment that
determines whether your current diffuse heatmap is caused by averaging or by the network. What
result would tell you the network genuinely has nothing sharp to say?

\exhead[value magnitude]{ch14:valuemag} Square $i$ attends $0.8$ to square $j$, but
$\lVert\vv{v}_j\rVert$ is a tenth of the typical value norm. How much of $i$'s new content
comes from $j$, relative to a square attended at $0.2$ with a typical value vector? What
should the heatmap have shown instead?

\exhead[sinks]{ch14:sinks} You notice that square h1 receives high attention in $70\%$ of the
positions you look at, across many heads. Give two hypotheses and the experiment that
separates them. Which would you bet on, and why?

\exhead[the diagonal]{ch14:diagonal} Should $\alpha_{ii}$ be included in a heatmap? Argue both
ways, then propose a display convention and justify it in one sentence.

\exhead[orientation]{ch14:orient} A tool shows attention for a black-to-move position and the
heatmap glows around White's king. Give the two possible explanations. Describe a two-minute
test using a single position that settles it.

\exhead[frames]{ch14:frames} List every frame-dependent quantity you have met in this book
(there are at least four). For each, name the convention and one symptom of getting it wrong.

\exhead[a contrast map]{ch14:contrast} Design a contrast experiment for the Morphy position
before 16.Qb8+: what modified position would you compare against, and what would you expect to
see if the network's attention is tracking the mating idea? What would you conclude if the two
maps were identical?

\exhead[an honest caption]{ch14:caption} Write the caption your frontend should display under
an attention heatmap. It must be short enough to fit and true enough to survive this chapter.

\exhead[the design question]{ch14:designq} Your project's design notes argue that a purely
visual translation may be more honest than a textual one, because a picture makes fewer
explicit claims. Using this chapter, argue the opposite: give two ways a picture can make a
\emph{stronger} false claim than a sentence would.

\end{exercises}
